\subsection{Funtores hipertor locales de complejos de módulos cuasicoherentes: caso general}
\label{subsection:III.6.5}

\begin{env}[6.5.1]
\label{III.6.5.1}
Sea ahora $S$ un preesquema arbitrario, y sean $X$ e $Y$ dos
$S$-preesquemas. Sean $\sh P_\bullet$ un complejo de
$\sh O_X$-módulos cuasicoherentes y $\sh Q_\bullet$ un complejo
de $\sh O_Y$-módulos cuasicoherentes. Pongamos $Z=X\times_SY$.
Vamos a definir $\sh O_Z$-módulos cuasicoherentes
\[
  \shTor_n^S(\sh P_\bullet,\sh Q_\bullet),
\]
denominados \emph{hipertor locales} de $\sh P_\bullet$ y
$\sh Q_\bullet$, que se reducen a los de \sref{III.6.4} cuando
$S$, $X$ e $Y$ son afines.

Cuando $\sh P_\bullet$ y $\sh Q_\bullet$ están concentrados en
grado cero, con términos respectivos $\sh F$ y $\sh G$, escribimos
$\shTor_n^S(\sh F,\sh G)$ en lugar de
$\shTor_n^S(\sh P_\bullet,\sh Q_\bullet)$.
\end{env}

\begin{env}[6.5.2]
\label{III.6.5.2}
Supongamos primero que $S$ es afín, y sean $(X_\lambda)$ y
$(Y_\mu)$ recubrimientos de $X$ e $Y$ por abiertos afines.
Entonces
\[
  Z_{\lambda\mu}=X_\lambda\times_SY_\mu
\]
forma un recubrimiento abierto afín de $Z$. Pongamos
\[
  \sh P_{\lambda,\bullet}=\sh P_\bullet|X_\lambda,
  \qquad
  \sh Q_{\mu,\bullet}=\sh Q_\bullet|Y_\mu.
\]
Para cada par $(\lambda,\mu)$ tenemos así un
$\sh O_{Z_{\lambda\mu}}$-módulo cuasicoherente
\[
  \sh F_{\lambda\mu}
  =
  \shTor_n^S(\sh P_{\lambda,\bullet},\sh Q_{\mu,\bullet}).
\]
Debemos mostrar que los $\sh F_{\lambda\mu}$ satisfacen la
condición de pegado \sref[0]{0.3.3.1}. Basta comprobar que, para
todo abierto afín
\[
  U\subset X_\lambda\cap X_{\lambda'},
  \qquad
  V\subset Y_\mu\cap Y_{\mu'},
\]
las restricciones de $\sh F_{\lambda\mu}$ y
$\sh F_{\lambda'\mu'}$ a $U\times_SV$ son canónicamente
isomorfas. Esto se deduce inmediatamente de sus isomorfismos
canónicos con
\[
  \shTor_n^S(\sh P_\bullet|U,\sh Q_\bullet|V),
\]
dados por \sref{III.6.4.5}.

La definición y \sref{III.6.4.5} muestran también que el
$\sh O_Z$-módulo así obtenido es independiente, salvo isomorfismo
único compatible con las identificaciones locales, de los
recubrimientos abiertos elegidos.
\oldpage[III]{149}
Lo denotamos por
\[
  \shTor_n^S(\sh P_\bullet,\sh Q_\bullet).
\]
Finalmente, para todo abierto $U$ de $X$ y todo abierto $V$ de
$Y$, su restricción a $U\times_SV$ es canónicamente isomorfa a
\[
  \shTor_n^S(\sh P_\bullet|U,\sh Q_\bullet|V).
\]
\end{env}

\begin{env}[6.5.3]
\label{III.6.5.3}
Pasemos ahora al caso general, en que $S$ es arbitrario. Sea
$(S_\alpha)$ un recubrimiento de $S$ por abiertos afines, y sean
$X_\alpha$ e $Y_\alpha$ las imágenes inversas de $S_\alpha$ en
$X$ e $Y$. Debemos mostrar que los haces
\[
  \sh G_\alpha
  =
  \shTor_n^{S_\alpha}
  (\sh P_\bullet|X_\alpha,\sh Q_\bullet|Y_\alpha)
\]
satisfacen la condición de pegado.

Para todo abierto afín $T\subset S_\alpha\cap S_\beta$, sean
$U$ y $V$ sus imágenes inversas en $X$ e $Y$. Basta definir
isomorfismos canónicos entre las restricciones de $\sh G_\alpha$
y $\sh G_\beta$ a $U\times_SV$ y
\[
  \shTor_n^T(\sh P_\bullet|U,\sh Q_\bullet|V).
\]
Podemos restringirnos además al caso en que
$T=\mathrm D(f_\alpha)=\mathrm D(f_\beta)$, donde $f_\alpha$ y
$f_\beta$ son secciones de $\sh O_S$ sobre $S_\alpha$ y
$S_\beta$, respectivamente. Por \sref{III.6.4.2},
\[
  \shTor_n^T(\sh P_\bullet|U,\sh Q_\bullet|V)
\]
es canónicamente isomorfo tanto a
$\shTor_n^{S_\alpha}(\sh P_\bullet|U,\sh Q_\bullet|V)$ como a
$\shTor_n^{S_\beta}(\sh P_\bullet|U,\sh Q_\bullet|V)$.
Combinando estos isomorfismos con las identificaciones locales de
\sref{III.6.5.2}, queda definido el $\sh O_Z$-módulo
\[
  \shTor_n^S(\sh P_\bullet,\sh Q_\bullet).
\]
Para todo abierto $U$ de $X$ y todo abierto $V$ de $Y$, la
restricción de este módulo a $U\times_SV$ es canónicamente
isomorfa a
\[
  \shTor_n^S(\sh P_\bullet|U,\sh Q_\bullet|V).
\]

Esta construcción define manifiestamente, en las categorías de
complejos acotados inferiormente de módulos cuasicoherentes, un
bi-$\partial$-funtor
\[
  \shTor_\bullet^S(\sh P_\bullet,\sh Q_\bullet)
\]
con valores en los $\sh O_Z$-módulos. En efecto, la afirmación es
local en $X$, $Y$ y $S$, por \sref{III.6.4.5} y por el hecho de
que \sref{III.6.4.4.2} es un morfismo de bi-$\partial$-funtores.

Si $\sh P_\bullet$ y $\sh Q_\bullet$ están concentrados en grado
cero, con términos respectivos $\sh F$ y $\sh G$, entonces
\[
  \shTor_0^S(\sh F,\sh G)
  =
  \sh F\otimes_S\sh G,
\]
el producto tensorial externo definido en \sref[I]{I.9.1.2}; esto
se deduce de \sref{III.6.4.1.1} y \sref[I]{I.9.1.3}.
\end{env}

\begin{env}[6.5.4]
\label{III.6.5.4}
La construcción precedente y \sref{III.6.4.6} muestran que
$\shTor_\bullet^S(\sh P_\bullet,\sh Q_\bullet)$ es el objeto al
que convergen dos funtores espectrales
$'\sh E(\sh P_\bullet,\sh Q_\bullet)$ y
$''\sh E(\sh P_\bullet,\sh Q_\bullet)$, cuyos términos $E_2$ son
\[
  '\sh E_{pq}^2
  =
  \sh H_p\bigl(\shTor_q^S(\sh P_\bullet,\sh Q_\bullet)\bigr)
  \tag{6.5.4.1}\label{III.6.5.4.1}
\]
y
\[
  ''\sh E_{pq}^2
  =
  \bigoplus_{q'+q''=q}
  \shTor_p^S
  \bigl(\sh H_{q'}(\sh P_\bullet),
        \sh H_{q''}(\sh Q_\bullet)\bigr).
  \tag{6.5.4.2}\label{III.6.5.4.2}
\]
La sucesión espectral \sref{III.6.5.4.2} es siempre regular.
Ambas sucesiones espectrales son birregulares cuando
$\sh P_\bullet$ y $\sh Q_\bullet$ están acotados inferiormente.
A continuación se describe otro caso en que son birregulares.
\end{env}

\begin{env}[6.5.5]
\label{III.6.5.5}
Sea $T$ un espacio topológico. Se dice que un haz de anillos
$\sh A$ sobre $T$ tiene \emph{dimensión cohomológica a lo sumo
$n$} si el anillo $\sh A_t$ tiene dimensión cohomológica a lo
sumo $n$ para todo $t\in T$. En este caso se dice también que el
espacio anillado $(T,\sh A)$ tiene dimensión cohomológica a lo
sumo $n$. Un haz de anillos, o un espacio anillado, tiene
\emph{dimensión cohomológica finita} si esto se cumple para algún
entero $n$.

Si los $\sh A_t$ son anillos locales noetherianos conmutativos,
\oldpage[III]{150}
decir que tienen dimensión cohomológica a lo sumo $n$ significa
que son regulares y tienen dimensión de Krull a lo sumo $n$
(\sref[0]{0.17.3.1}). En la terminología de la teoría de la
dimensión que se introducirá en el capítulo IV, un preesquema
localmente noetheriano $T$ tiene dimensión cohomológica a lo sumo
$n$ exactamente cuando es regular (\sref[0]{0.4.1.4}) y tiene
dimensión a lo sumo $n$. Equivalentemente, para todo abierto afín
$U$ de $T$, el anillo $\Gamma(U,\sh O_T)$ tiene dimensión
cohomológica a lo sumo $n$ (\sref[0]{0.17.2.6}).

En consecuencia, si $S$ es localmente noetheriano y tiene
dimensión cohomológica finita, \sref{III.6.3.2} muestra que tanto
$'\sh E(\sh P_\bullet,\sh Q_\bullet)$ como
$''\sh E(\sh P_\bullet,\sh Q_\bullet)$ son birregulares.

Finalmente, bajo el isomorfismo canónico
$X\times_SY\xrightarrow{\sim}Y\times_SX$,
$\shTor_\bullet^S(\sh P_\bullet,\sh Q_\bullet)$ corresponde
canónicamente a
$\shTor_\bullet^S(\sh Q_\bullet,\sh P_\bullet)$.
\end{env}

\begin{proposition}[6.5.6]
\label{III.6.5.6}
Sea $(\sh P_{\alpha\bullet})$ un sistema inductivo filtrante de
complejos de $\sh O_X$-módulos cuasicoherentes. Existe un
isomorfismo canónico
\[
  \varinjlim_\alpha
  \shTor_\bullet^S(\sh P_{\alpha\bullet},\sh Q_\bullet)
  \xrightarrow{\sim}
  \shTor_\bullet^S
  \left(\varinjlim_\alpha\sh P_{\alpha\bullet},\sh Q_\bullet\right).
  \tag{6.5.6.1}\label{III.6.5.6.1}
\]
\end{proposition}

\begin{proof}
La cuestión es local en $S$, $X$ e $Y$. Por tanto, podemos
suponer que $S$, $X$ e $Y$ son afines, en cuyo caso la proposición
se reduce a \sref{III.6.3.6}.
\end{proof}

\begin{env}[6.5.7]
\label{III.6.5.7}
\emph{Observaciones.}

\textnormal{(i)} Consideremos en particular el caso $S=X=Y$, con
$\sh P_\bullet$ y $\sh Q_\bullet$ dos complejos de
$\sh O_S$-módulos cuasicoherentes. Entonces los
$\shTor_n^S(\sh P_\bullet,\sh Q_\bullet)$ son
$\sh O_S$-módulos cuasicoherentes. Además, para todo punto
$z\in S$, \sref{III.6.5.6} da un isomorfismo canónico
\[
  \bigl(\shTor_n^S(\sh P_\bullet,\sh Q_\bullet)\bigr)_z
  \xrightarrow{\sim}
  \mathbf{Tor}_n^{\sh O_z}
  \bigl((\sh P_\bullet)_z,(\sh Q_\bullet)_z\bigr),
  \tag{6.5.7.1}\label{III.6.5.7.1}
\]
porque la cuestión es local y, por \sref{III.6.4.1.1}, se reduce
al caso de módulos.

\textnormal{(ii)} La definición de hipertor puede generalizarse
a dos complejos $\sh P_\bullet,\sh Q_\bullet$ de
$\sh O_X$-módulos sobre un mismo espacio anillado $(X,\sh O_X)$.
Para cada abierto $U$ de $X$, pongamos
\[
  A(U)=\Gamma(U,\sh O_X),\qquad
  P_\bullet(U)=\Gamma(U,\sh P_\bullet),\qquad
  Q_\bullet(U)=\Gamma(U,\sh Q_\bullet).
\]
Los $A(U)$-módulos
\[
  \mathbf{Tor}_n^{A(U)}\bigl(P_\bullet(U),Q_\bullet(U)\bigr)
\]
forman entonces un prehaz sobre $X$; denotemos por
$\shTor_n^X(\sh P_\bullet,\sh Q_\bullet)$ el $\sh O_X$-módulo
asociado. Cuando $X$ es un preesquema, \sref{III.6.3.12} muestra
que este módulo es canónicamente isomorfo al hipertor definido
más arriba. No desarrollaremos más esta generalización.
\end{env}

\begin{proposition}[6.5.8]
\label{III.6.5.8}
Sean $X$ e $Y$ dos $S$-preesquemas, y sea $\sh F$
(respectivamente $\sh G$) un $\sh O_X$-módulo cuasicoherente
(respectivamente un $\sh O_Y$-módulo cuasicoherente). Si
$\sh F$ o $\sh G$ es plano sobre $S$, entonces
\[
  \shTor_n^S(\sh F,\sh G)=0\qquad(n\ne0).
\]
\end{proposition}

\begin{proof}
La cuestión es local en $X$ e $Y$, por lo que podemos suponer que
$X$, $Y$ y $S$ son afines, de anillos respectivos $B$, $C$ y
$A$, y escribir $\sh F=\widetilde M$ y $\sh G=\widetilde N$,
donde $M$ es un $B$-módulo y $N$ un $C$-módulo. Supongamos, por
ejemplo, que $\sh F$ es plano sobre $S$. Esto significa que, para
todo $s\in S$, $M_s$ es un $A_s$-módulo plano
(\sref[0]{0.6.7.1}); por consiguiente, $M$ es un $A$-módulo plano
(\sref[0]{0.6.3.3}). Así,
$\Tor_n^A(M,N)=0$ para $n>0$ y todo $C$-módulo $N$
(\sref[0]{0.6.1.1}), y la conclusión se deduce de
\sref{III.6.4.1.1}.
\end{proof}

\begin{corollary}[6.5.9]
\label{III.6.5.9}
Sean $X$ e $Y$ dos $S$-preesquemas, y sea $\sh P_\bullet$
(respectivamente $\sh Q_\bullet$) un complejo
\oldpage[III]{151}
acotado
inferiormente de $\sh O_X$-módulos cuasicoherentes
(respectivamente de $\sh O_Y$-módulos cuasicoherentes).
Supongamos que cada $\sh P_i$ es plano sobre $S$. Entonces existe
un isomorfismo canónico de $\partial$-funtores en $\sh Q_\bullet$
\[
  \shTor_\bullet^S(\sh P_\bullet,\sh Q_\bullet)
  \xrightarrow{\sim}
  \sh H_\bullet(\sh P_\bullet\otimes_S\sh Q_\bullet).
  \tag{6.5.9.1}\label{III.6.5.9.1}
\]
\end{corollary}

\begin{proof}
Cuando $S$, $X$ e $Y$ son afines, esto es precisamente
\sref{III.6.3.7}; el caso general se deduce mediante los
razonamientos de \sref{III.6.5.2} y \sref{III.6.5.3}.
\end{proof}

\begin{env}[6.5.10]
\label{III.6.5.10}
\emph{Corolario.}
Supongamos que $X$ es plano sobre $S$ (\sref[0]{0.6.7.1}), que
$\sh P_\bullet$ y $\sh Q_\bullet$ están acotados inferiormente,
y que cada $\sh P_i$ es un $\sh O_X$-módulo localmente libre, no
necesariamente de tipo finito. Entonces el homomorfismo
\sref{III.6.5.9.1} es biyectivo.

En efecto, se cumple la hipótesis de \sref{III.6.5.9}: la
planitud es, por definición, una propiedad puntual en $X$, y toda
suma directa de módulos planos es plana (\sref[0]{0.6.1.2}).
\end{env}

\begin{proposition}[6.5.11]
\label{III.6.5.11}
Sean $X'$ e $Y'$ dos $S$-preesquemas, y sean
$f:X\to X'$ y $g:Y\to Y'$ dos $S$-morfismos afines. Para un
complejo $\sh P_\bullet$ de $\sh O_X$-módulos cuasicoherentes y
un complejo $\sh Q_\bullet$ de $\sh O_Y$-módulos
cuasicoherentes, existe un isomorfismo canónico funtorial
\[
  (f\times_Sg)_*
  \bigl(\shTor_\bullet^S(\sh P_\bullet,\sh Q_\bullet)\bigr)
  \xrightarrow{\sim}
  \shTor_\bullet^S(f_*\sh P_\bullet,g_*\sh Q_\bullet).
  \tag{6.5.11.1}\label{III.6.5.11.1}
\]
\end{proposition}

\begin{proof}
Como $f$ y $g$ son afines, $f_*\sh P_\bullet$ y
$g_*\sh Q_\bullet$ son complejos de módulos cuasicoherentes
(\sref[II]{II.1.2.6}). Si $Z'=X'\times_SY'$, ambos miembros de
\sref{III.6.5.11.1} son $\sh O_{Z'}$-módulos cuasicoherentes por
\sref{III.6.5.1}. Nos reducimos inmediatamente al caso en que
$S$, $X'$ e $Y'$ son afines; entonces $X$ e $Y$ también son
afines, y la afirmación se deduce de \sref{III.6.4.1.1} y
\sref[I]{I.1.6.3}.
\end{proof}

\begin{env}[6.5.12]
\label{III.6.5.12}
\emph{Observación.}
Sean $X'$ e $Y'$ dos $S$-preesquemas, con $X'$ plano sobre $S$.
Sea $X$ un subpreesquema cerrado de $X'$, y sea
$i:X\to X'$ la inmersión canónica. Sea $\sh P_\bullet$
(respectivamente $\sh Q_\bullet$) un complejo acotado
inferiormente de $\sh O_X$-módulos cuasicoherentes
(respectivamente de $\sh O_{Y'}$-módulos cuasicoherentes).
Finalmente, sea $\sh L'_{\bullet\bullet}$ una resolución de
$i_*\sh P_\bullet$ formada por $\sh O_{X'}$-módulos localmente
libres, tal que cada punto de $X'$ tenga un entorno abierto afín
$U$ para el cual $\sh L'_{j,\bullet}|U$ sea una resolución libre de
$i_*\sh P_j|U$, para todo $j$. Existe entonces un isomorfismo
canónico
\[
  (i\times_S1)_*
  \bigl(\shTor_\bullet^S(\sh P_\bullet,\sh Q_\bullet)\bigr)
  \xrightarrow{\sim}
  \sh H_\bullet(\sh L'_{\bullet\bullet}\otimes_S\sh Q_\bullet).
  \tag{6.5.12.1}\label{III.6.5.12.1}
\]

Si $S$, $X'$ e $Y'$ son afines, y si
$\sh L'_{j,\bullet}$ es una resolución libre de $i_*\sh P_j$ para
todo $j$, \sref{III.6.5.11} nos reduce al caso $X'=X$, con $X$
plano sobre $S$, y basta aplicar \sref{III.6.3.4}. En general,
\sref{III.6.5.12.1} se define localmente, y queda comprobar que
estas definiciones locales dan un isomorfismo global.

Para ello, recordemos la definición del primer isomorfismo
\sref{III.6.3.4.1}. Proviene de un isomorfismo de sucesiones
espectrales (\sref[0]{0.11.6.5} y \sref[0]{0.11.5.3}), obtenido a
su vez de un morfismo de bicomplejos
\[
  L_{\bullet\bullet}\longrightarrow L''_{\bullet\bullet},
\]
donde $L_{\bullet\bullet}$ es la resolución dada de $P_\bullet$,
y $L''_{\bullet\bullet}$ es una resolución proyectiva de
$P_\bullet$ en la categoría de complejos acotados
\oldpage[III]{152}
inferiormente
de $A$-módulos (\sref[0]{0.11.5.2.2}). La afirmación se deduce del
hecho de que \sref{III.6.3.4.1} es independiente de la resolución
proyectiva elegida $L''_{\bullet\bullet}$, pues existe una equivalencia
homotópica entre cualesquiera dos resoluciones de este tipo
(\emph{M}, V, 1.2).
\end{env}

\begin{proposition}[6.5.13]
\label{III.6.5.13}
Sean $X$ e $Y$ dos $S$-preesquemas, y supongamos que se cumple
una de las condiciones siguientes:
\begin{enumerate}
  \item[(i)] $X$ y $Z=X\times_SY$ son localmente noetherianos, y
    $X$ es plano sobre $S$;
  \item[(ii)] $S$ y $X$ son localmente noetherianos, y $Y$ es de
    tipo finito sobre $S$.
\end{enumerate}
Sea $\sh P_\bullet$ (respectivamente $\sh Q_\bullet$) un complejo
acotado inferiormente de $\sh O_X$-módulos cuasicoherentes
(respectivamente de $\sh O_Y$-módulos cuasicoherentes).
Supongamos además que, para todo $n$,
$\sh H_n(\sh P_\bullet)$ (respectivamente
$\sh H_n(\sh Q_\bullet)$) es un $\sh O_X$-módulo
(respectivamente un $\sh O_Y$-módulo) de tipo finito. Entonces
\[
  \shTor_n^S(\sh P_\bullet,\sh Q_\bullet)
\]
es un $\sh O_Z$-módulo coherente para todo $n$.
\end{proposition}

\begin{proof}
Como $\sh P_\bullet$ y $\sh Q_\bullet$ están acotados
inferiormente, la sucesión espectral
$''\sh E(\sh P_\bullet,\sh Q_\bullet)$ es birregular por
\sref{III.6.5.4}. Tanto en el caso (i) como en el caso (ii), $Z$
es localmente noetheriano, de modo que \sref[0]{0.11.1.8} reduce
la afirmación a la coherencia de los términos
$''\sh E_{pq}^2$. La hipótesis de tipo finito sobre los módulos de
homología y la fórmula \sref{III.6.5.4.2} reducen la proposición
al caso particular en que $\sh P_\bullet$ y $\sh Q_\bullet$ están
concentrados en grado cero, es decir, al corolario siguiente.
\end{proof}

\begin{corollary}[6.5.14]
\label{III.6.5.14}
Supongamos que se cumple una de las condiciones \emph{(i)} y
\emph{(ii)} de \sref{III.6.5.13}, y sea $\sh F$
(respectivamente $\sh G$) un $\sh O_X$-módulo cuasicoherente
(respectivamente un $\sh O_Y$-módulo cuasicoherente) de tipo
finito. Entonces $\shTor_n^S(\sh F,\sh G)$ es un
$\sh O_Z$-módulo coherente para todo $n$.
\end{corollary}

\begin{proof}
La cuestión es local en $X$ e $Y$, por lo que podemos suponer que
$S$, $X$ e $Y$ son afines.

\textnormal{(i)} Bajo las hipótesis de (i), $S$, $X$ y $Z$ son
noetherianos. Existe, por tanto, una resolución localmente libre
$\sh L_\bullet$ de $\sh F$, formada por $\sh O_X$-módulos de
tipo finito (\sref[I]{I.1.3.7}). Como $X$ es plano sobre $S$,
\sref{III.6.3.4} da
\[
  \shTor_n^S(\sh F,\sh G)
  =
  \sh H_n(\sh L_\bullet\otimes_S\sh G).
\]
Los módulos $\sh L_i\otimes_S\sh G$ son
$\sh O_Z$-módulos cuasicoherentes de tipo finito
(\sref[I]{I.9.1.1}), luego son coherentes. Se sigue que
$\sh H_n(\sh L_\bullet\otimes_S\sh G)$ es coherente
(\sref[0]{0.5.3.4}).

\textnormal{(ii)} Supongamos ahora que se cumple la condición
(ii). Como $A(Y)$ es cociente de una $A(S)$-álgebra de polinomios
en un número finito de indeterminadas (\sref[I]{I.6.3.3}), $Y$ es
un subpreesquema cerrado de un $S$-esquema afín $Y'$ plano y de
tipo finito sobre $S$. El esquema $Y'$ es noetheriano
(\sref[I]{I.6.3.7}), por lo que existe una resolución localmente
libre $\sh M_\bullet$ de $j_*\sh G$ por
$\sh O_{Y'}$-módulos de tipo finito, donde $j:Y\to Y'$ es la
inmersión canónica. Por \sref{III.6.5.12},
$\shTor_n^S(\sh F,\sh G)$ es la imagen inversa por $1\times j$
del $\sh O_{Z'}$-módulo
\[
  \sh H_n(\sh F\otimes_S\sh M_\bullet),
  \qquad Z'=X\times_SY'.
\]
Como en (i), los $\sh F\otimes_S\sh M_i$ son
$\sh O_{Z'}$-módulos coherentes, y \sref[0]{0.5.3.4} da de nuevo
la conclusión.
\end{proof}

\begin{env}[6.5.15]
\label{III.6.5.15}
La teoría desarrollada arriba para dos complejos
$\sh P_\bullet,\sh Q_\bullet$ de haces cuasicoherentes sobre dos
$S$-preesquemas $X$ e $Y$ se extiende sin dificultad a una
familia finita arbitraria de $S$-preesquemas $X^{(i)}$
$(1\leq i\leq m)$, cada uno provisto de un complejo
$\sh P_\bullet^{(i)}$ de $\sh O_{X^{(i)}}$-módulos
cuasicoherentes. Si
\[
  Z=X^{(1)}\times_SX^{(2)}\times_S\cdots\times_SX^{(m)},
\]
se define así un $\sh O_Z$-módulo cuasicoherente
\[
  \shTor_q^S
  \bigl(\sh P_\bullet^{(1)},\sh P_\bullet^{(2)},\ldots,
        \sh P_\bullet^{(m)}\bigr).
\]
Dejamos al lector el desarrollo de la teoría en este caso general,
y registramos únicamente, para referencia posterior, el término
$E_2$ de la segunda sucesión espectral regular que converge a este
hipertor:

\oldpage[III]{153}

\[
  ''\sh E_{pq}^2
  =
  \bigoplus_{q_1+q_2+\cdots+q_m=q}
  \shTor_p^S
  \bigl(
    \sh H_{q_1}(\sh P_\bullet^{(1)}),\ldots,
    \sh H_{q_m}(\sh P_\bullet^{(m)})
  \bigr).
  \tag{6.5.15.1}\label{III.6.5.15.1}
\]
En \sref{III.6.8} estudiaremos las sucesiones espectrales de
asociatividad a las que dan lugar estos funtores hipertor de un
número finito arbitrario de complejos.
\end{env}
